\chapter{Analyse et spécification des besoins}

\section{Acteurs du système}

\noindent
La plateforme \projectName{} repose sur trois rôles principaux.

\begin{table}[h!]
\centering
\renewcommand{\arraystretch}{1.4}
\begin{tabular}{|p{3cm}|p{10cm}|}
\hline
\textbf{Acteur} & \textbf{Description} \\
\hline
Candidat & Utilisateur qui cherche une opportunité dans le secteur CHR. Il complète son profil, ajoute son CV, consulte les offres et postule. \\
\hline
Recruteur & Établissement ou responsable RH qui publie des offres, consulte les candidatures et prend des décisions. \\
\hline
Administrateur & Responsable de la supervision globale, de la modération des offres et du suivi des utilisateurs. \\
\hline
\end{tabular}
\caption{Acteurs principaux de la plateforme}
\label{tab:acteurs}
\end{table}

\section{Besoins fonctionnels}

\subsection{Espace candidat}

\begin{itemize}
    \item créer un compte et se connecter ;
    \item compléter un profil professionnel : ville, expérience, poste recherché, disponibilité et contrat préféré ;
    \item ajouter un CV PDF et une photo optionnelle ;
    \item consulter, filtrer et sauvegarder les offres ;
    \item postuler à une offre en réutilisant le CV du profil ;
    \item passer un quiz lorsque l'offre le demande ;
    \item consulter les candidatures, les statuts et les notifications.
\end{itemize}

\subsection{Espace recruteur}

\begin{itemize}
    \item compléter le profil de l'établissement : nom, ville et type ;
    \item créer, modifier et suivre les offres publiées ;
    \item ajouter une image optionnelle à une offre ;
    \item créer un quiz QCM optionnel ;
    \item consulter les candidatures reçues, les CV et les scores de quiz ;
    \item accepter ou refuser une candidature ;
    \item consulter le quota de vues et l'offre premium.
\end{itemize}

\subsection{Espace administrateur}

\begin{itemize}
    \item consulter les statistiques globales ;
    \item filtrer et rechercher les utilisateurs ;
    \item modérer les offres ;
    \item suspendre une offre avec un motif optionnel ;
    \item réactiver une offre suspendue.
\end{itemize}

\section{Besoins non fonctionnels}

\noindent
Les besoins non fonctionnels définissent les qualités attendues du système :

\begin{itemize}
    \item \textbf{Sécurité :} accès aux pages et API selon le rôle de l'utilisateur ;
    \item \textbf{Fiabilité :} protection contre les candidatures dupliquées ;
    \item \textbf{Ergonomie :} interface claire et responsive ;
    \item \textbf{Performance :} chargement rapide des pages et requêtes API simples ;
    \item \textbf{Maintenabilité :} séparation claire entre frontend, backend et base de données ;
    \item \textbf{Accessibilité de démonstration :} parcours compréhensible pour la soutenance.
\end{itemize}

\section{Règles métier principales}

\begin{itemize}
    \item un candidat doit compléter son profil et ajouter un CV avant de postuler ;
    \item le CV est ajouté une seule fois dans le profil et réutilisé dans les candidatures ;
    \item un candidat ne peut pas postuler deux fois à la même offre ;
    \item le quiz est passé après la création de la candidature ;
    \item les offres suspendues ou expirées ne doivent pas apparaître dans la liste publique active ;
    \item un recruteur gratuit est limité dans la consultation des profils candidats ;
    \item la discussion n'est disponible qu'après acceptation d'une candidature.
\end{itemize}

\section{Cas d'utilisation principaux}

\begin{table}[h!]
\centering
\renewcommand{\arraystretch}{1.4}
\begin{tabular}{|p{4cm}|p{9cm}|}
\hline
\textbf{Cas d'utilisation} & \textbf{Description} \\
\hline
Postuler à une offre & Le candidat consulte une offre, vérifie son profil et confirme la postulation. \\
\hline
Passer un quiz & Le candidat répond aux questions QCM associées à une offre après avoir postulé. \\
\hline
Créer une offre & Le recruteur renseigne les détails de l'offre et peut ajouter un quiz. \\
\hline
Traiter une candidature & Le recruteur consulte le CV, le score du quiz et accepte ou refuse le candidat. \\
\hline
Modérer une offre & L'administrateur suspend ou active une offre selon son état. \\
\hline
\end{tabular}
\caption{Cas d'utilisation principaux}
\label{tab:cas_utilisation}
\end{table}
